\subsection{拓扑群中点的邻域}

设 \(\mathfrak{B}\) 为拓扑群 G 中单位元 \(e\) 的邻域滤子，并设 \(a\) 为 G 的任意一点。由于 \(x \to ax\) 与 \(x \to xa\) 都是同胚，因此 \(a\) 的邻域滤子是族 \(a\) . \(\mathfrak{B}\)，即由诸集合 \(a\) . V（其中 V 取遍 \(\mathfrak{B}\)）组成，它同时也是族 \(\mathfrak{B}\) . \(a\)，即由诸集合 V. \(a\) 组成。由此可见，只要知道了群的单位元 \(e\) 的邻域滤子，我们便知道了拓扑群中任意一点的邻域滤子。

如果我们说 \(xy\) 与 \(x^{-1}\) 在 \(x = y = e\) 处连续，就得到（第一章 § 2 第 1 目）：

\((\mathbf{GV}_{\mathbf{I}})\). 对任意 \(\mathbf{U} \in \mathfrak{B}\)，存在 \(\mathbf{V} \in \mathfrak{B}\) 使 \(\mathbf{V} \cdot \mathbf{V} \subset \mathbf{U}\)。

\((\mathbf{GV}_{\mathbf{II}})\). 对任意 \(\mathbf{U}\in \mathfrak{B},\) 有 \(\mathbf{U}^{-1}\in \mathfrak{B}\)

任何滤子 \(\mathfrak{B}\)，只要它在 \(G\) 上并且满足 \((\mathrm{GV}_{\mathbb{I}})\) 与 \((\mathrm{GV}_{\mathbb{II}})\)，便也满足 \((\mathrm{GV}_{\alpha})\)：对任意 \(U \in \mathfrak{B}\)，存在 \(V \in \mathfrak{B}\) 使 \(V \cdot V^{-1} \subset U\)。因为由 \((\mathrm{GV}_{\mathbb{I}})\)，存在 \(W \in \mathfrak{B}\) 使 \(W \cdot W \subset U\)；又由 \((\mathrm{GV}_{\mathbb{II}})\)，存在 \(V \in \mathfrak{B}\) 使 \(V \subset W \cap W^{-1}\)；于是 \(V^{-1} \subset W\)，从而 \(V \cdot V^{-1} \subset W \cdot W \subset U\)。

反之，设 \(\mathfrak{B}\) 为 \(\mathbf{G}\) 上满足 \((\mathrm{GV}_{\alpha})\) 的一个滤子，则首先可知 \(e\) 属于每个集合 \(\mathbf{U} \in \mathfrak{B}\)；因为若 \(\mathbf{V} \in \mathfrak{B}\) 使 \(\mathbf{V} \cdot \mathbf{V}^{-1} \subset \mathbf{U}\)，则由于 V 非空，我们有 \(x.x^{-1}=e\in U\)，它对每个 \(x\in V\) 成立。于是条件 \((\mathrm{GV}_{\alpha})\) 蕴含 \(V^{-1}\subset V.V^{-1}\subset U\)，因而 \(U^{-1}\in\mathfrak{B}\)（只要 \(U\in\mathfrak{B}\)）。最后，若 \(V\in\mathfrak{B}\) 使 \(V.V^{-1}\subset U\)，且 \(W\in\mathfrak{B}\) 使 \(W\subset V\cap V^{-1}\)，则 \(W.W\subset U\)。由此可见，\((\mathrm{GV}_{\alpha})\) 等价于 \((\mathrm{GV}_{I})\) 与 \((\mathrm{GV}_{II})\) 的合取。

最后，由于 \(x \rightarrow axa^{-1}\) 是使 e 固定不动的同胚，B 具有下列性质：

\((\mathrm{GV}_{\mathrm{III}})\). 对所有 \(a \in \mathbf{G}\) 与所有 \(\mathbf{V} \in \mathfrak{B}\)，有 \(a. \mathbf{V}.a^{-1} \in \mathfrak{B}\)。

滤子 \(\mathfrak{B}\) 的这三个性质是特征性质：

命题 1 设 G 是一个群，B 是 G 上满足公理 \((\mathrm{GV}_{\mathrm{I}})\)、\((\mathrm{GV}_{\mathrm{II}})\) 与 \((\mathrm{GV}_{\mathrm{III}})\) 的一个滤子。则在 G 上存在唯一的与 G 的群结构相容的拓扑，使 B 是单位元 e 的邻域滤子。对这个拓扑，任一点 \(a \in G\) 的邻域滤子与两个滤子 a.B 和 B.a 中的每一个都相同。

如果存在具有所需性质的拓扑，那么由上文所述，\(a\) 的邻域滤子与两个滤子 \(a\) . \(\mathfrak{B}\) 和 \(\mathfrak{B}.a\) 中的每一个都重合；因此，如果这样的拓扑存在，它便是唯一的。只要证明 1) 诸滤子 \(a\) . \(\mathfrak{B}\) 是 \(\mathbf{G}\) 上某个拓扑的邻域滤子，以及 2) 这个拓扑与 \(\mathbf{G}\) 的群结构相容，存在性便得证。

\begin{enumerate}
\def\labelenumi{\arabic{enumi})}
\item
  滤子 \(a. \mathfrak{B}\) 满足公理 \((\mathrm{V}_{\mathrm{III}})\)（见第一章 § 1 第 2 目），这由 \((\mathrm{GV}_{\mathrm{I}})\) 与 \((\mathrm{GV}_{\mathrm{II}})\) 推出，我们前面已经看到；因此，要证明 \(a. \mathfrak{B}\) 是 \(a\) 在 \(G\) 上一个拓扑中的邻域滤子，只需验证公理 \((\mathrm{V}_{\mathrm{IV}})\)。设 \(V\) 是 \(\mathfrak{B}\) 中任一集合，\(W\) 是 \(\mathfrak{B}\) 中使 \(W.W \subset V\) 的一个集合；则对任意 \(x \in a.W\)，有 \(x.W \subset a.W.W \subset a.V\)，于是 \(a.V\) 属于滤子 \(x. \mathfrak{B}\)；因而 \((\mathrm{V}_{\mathrm{IV}})\) 得到满足。
\item
  我们现在证明由邻域滤子 \(a\) . \(\mathfrak{B}\) 所定义的拓扑满足 (GT')。设 \(a, b\) 是 G 的任意两点；若令 \(x = au\)、\(y = bv\)，则我们要证明：\(xy^{-1}\) 可以随意地接近 \(ab^{-1}\)，只要 \(u\) 与 \(v\) 充分接近 \(e\)。现在 \((ab^{-1})^{-1}(xy^{-1}) = buv^{-1}b^{-1}\)；设 U 是 \(e\) 的任一邻域，则我们将有 \(buv^{-1}b^{-1} \in U\)，只要 \(uv^{-1} \in b^{-1}Ub = V\)，并且 \(V \in \mathfrak{B}\)（由 \((\mathrm{GV}_{\mathrm{III}})\)）。但由 \((\mathrm{GV}_{\mathrm{I}})\) 与 \((\mathrm{GV}_{\mathrm{II}})\)，存在 \(W \in \mathfrak{B}\) 使 \(W.W^{-1} \subset V\)；因此只需取 \(u \in W\) 与 \(v \in W\)，便有 \(xy^{-1} \in (ab^{-1})U\)。证毕。
\end{enumerate}

在 G 上定义一个与群结构相容的拓扑，一个常用方法是给出一个满足公理 \((\mathrm{GV}_{\mathrm{I}})\)、\((\mathrm{GV}_{\mathrm{II}})\) 与 \((\mathrm{GV}_{\mathrm{III}})\) 的滤子。对于滤子基 B，相应的条件如下：

\((\mathbf{GV}_{\mathbf{I}}^{\prime})\). 对任意 \(\mathbf{U}\in \mathfrak{B}\)，存在 \(\mathbf{V}\in \mathfrak{B}\) 使 \(\mathbf{V}.\mathbf{V}\subset \mathbf{U}\)。\((\mathrm{GV}_{\mathrm{II}}^{\prime})\). 对任意 \(\mathbf{U} \in \mathfrak{B}\)，存在 \(\mathbf{V} \in \mathfrak{B}\) 使 \(\mathbf{V}^{-1} \subset \mathbf{U}\)。\((\mathrm{GV}_{\mathrm{III}}^{\prime})\). 对任意 \(a \in \mathbf{G}\) 与任意 \(\mathbf{U} \in \mathfrak{B}\)，存在 \(\mathbf{V} \in \mathfrak{B}\) 使 \(\mathbf{V} \subset a. \mathbf{U}. a^{-1}\)。

e 的一个邻域若与它在对称映射 \(x \rightarrow x^{-1}\) 下的像重合，就称为对称邻域。若 V 是 e 的任一邻域，则 \(V \cup V^{-1}\)、\(V \cap V^{-1}\) 与 \(V \cdot V^{-1}\) 都是对称邻域。由 \((\mathrm{GV}_{\mathrm{II}})\)，诸对称邻域构成 e 的一个基本邻域系。又由 \((\mathrm{GV}_{\mathrm{I}})\)，当 V 取遍 e 的一个基本邻域系时，诸集合 \(V^{n}\)（其中 n 是一个固定的整数 \(\neq 0\)）构成 e 的一个基本邻域系。

注。若 G 是交换群，则有 \(x \cdot A \cdot x^{-1} = A\)（对 G 的每个子集 A 及每个 \(x \in G\)），因而 \((GV_{III})\){[}相应地 \((GV_{III}')\){]}对 G 上每个滤子（相应地，滤子基）自动满足。另一方面，若 G 不是交换群，则 \((GV_{III})\) 不是 \((GV_I)\) 与 \((GV_{II})\) 的推论{[}参看习题 5{]}。

若 G 是以加法记写的交换群，则刻画与 G 的群结构相容的拓扑的原点邻域滤子 B 的公理就是下列两条：\((\mathrm{GA}_{\mathrm{I}})\). 对任意 \(U \in B\)，存在 \(V \in B\) 使 \(V + V \subset U\)。\((\mathrm{GA}_{\mathrm{II}})\). 对任意 \(U \in B\)，有 \(-U \in B\)。

命题 2 拓扑群 G 是 Hausdorff 群，当且仅当集合 \(\{e\}\) 是闭集。

显然，若 G 是 Hausdorff 群，则 \(\{e\}\) 是闭集。反之，若 \(\{e\}\) 是闭集，则对角线 \(\Delta\)（\(G \times G\) 的）是闭集，因为它是 \(\{e\}\) 在连续映射 \((x, y) \to xy^{-1}\) 下的逆像；于是（第一章 § 8 第 1 目命题 1）G 是 Hausdorff 群。

推论 拓扑群 G 是 Hausdorff 群，当且仅当 e 的诸邻域之交仅由点 e 组成。

该条件显然必要。反之，若 \(e\) 的所有邻域之交恰为 \(\{e\}\)，则对任意 \(x \neq e\)，存在一个邻域 \(V\)（\(e\) 的），使 \(x^{-1} \notin V\)，因而 \(e \notin xV\)。这表明 \(x\) 不属于 \(\{e\}\) 的闭包，从而 \(\{e\}\) 是闭集，于是 \(G\) 是 Hausdorff 群。

例。借助一组子群来定义群上的拓扑。设 \(\mathfrak{B}\) 是群 \(\mathbf{G}\) 上由 \(\mathbf{G}\) 的子群组成的一个滤子基，则直接可以看出 \(\mathfrak{B}\) 满足公理 \((\mathrm{GV}_{\mathrm{I}}^{\prime})\) 与 \((\mathrm{GV}_{\mathrm{II}}^{\prime})\)，因为 \(\mathbf{H}.\mathbf{H}^{-1} = \mathbf{H}\)，其中 \(\mathbf{H}\) 是 \(\mathbf{G}\) 的任一子群。因此，集合 \(\mathfrak{B}\) 将是 \(e\) 在某个与 \(\mathbf{G}\) 的群结构相容的拓扑中的基本邻域系，只要 \(\mathfrak{B}\) 满足 \((\mathrm{GV}_{\mathrm{III}})\)；特别地，当 \(\mathfrak{B}\) 中所有子群都是正规子群时即属此情形，因而在 \(G\) 是交换群时永远如此。这样定义的拓扑，按命题 2，是 Hausdorff 的当且仅当 \(\mathfrak{B}\) 中所有子群之交仅由 \(e\) 组成。最有意义的情形是子群 \(\{e\}\) 不属于 \(\mathfrak{B}\) 的情形（否则由 \(\mathfrak{B}\) 定义的拓扑是离散拓扑）：若 \(\{e\} \notin \mathfrak{B}\)，则由 \(\mathfrak{B}\) 定义的拓扑只有在 \(\mathfrak{B}\) 是无限集时才可能是 Hausdorff 的。

由于两个子群之交是子群，我们可以从 G 的任意一组子群 \(\mathfrak{F}\) 出发，在 G 上定义一个与其群结构相容的拓扑：设 \(\mathfrak{G}\) 是所有形如 \(a.H.a^{-1}\) 的子群组成的集，其中 \(H\in\mathfrak{F}\) 且 \(a\in G\)；再设 \(\mathfrak{B}\) 是属于 \(\mathfrak{G}\) 的子群的所有有限交组成的集。则 \(\mathfrak{B}\) 是一个滤子基且满足 \((\mathrm{GV}_{\mathrm{III}}^{\prime})\)。

特别地，考虑环 A 的加法群。A 的理想的任意一个集合 \(\mathfrak{F}\) 都定义一个与这个加法群结构相容的拓扑。若 \(\mathfrak{F}\) 中所有理想之交是零理想，则这个拓扑是 Hausdorff 的；若 \(\mathfrak{F}\) 中诸理想的任何有限交都不是零理想，则它不是离散的。用这种方式定义的拓扑在数论中起很大作用（见本章 § 6 与 § 7 的习题）。

